%% ============================================================================
%%  applications/template.tex  —  SKELETON RESUME (single source template)
%% ----------------------------------------------------------------------------
%%  This is the ONLY resume template. The orchestrator copies this file into a
%%  position folder, then the writer fills the Experience / Projects / Skills
%%  sections from the context.md bullet pool. It replaces the old per-track
%%  templates (full-stack / ai-ml / dev-ops) — track is now a doctrine-routing
%%  label, not a file choice, so any track (including robotics/autonomy) fills
%%  this same skeleton.
%%
%%  WHAT IS FIXED (do not change): header (\name, \contact) and Education.
%%  WHAT GETS FILLED: Experience entries, Projects entries, Technical Skills.
%%
%%  rSectionEntry signature — 5 positional params, ALWAYS supply all 5 braces:
%%     {1 name}{2 dates}{3 tagline/role/link}{4 technologies}{5 location}
%%     - param 4 (technologies) renders after the name with a " | " separator;
%%       leave it {} for experiences, fill it for projects (the tech-stack line).
%%     - any unused param is an empty {}.
%%  rSet signature — {1 bucket label}{2 comma-separated items}.
%%
%%  PATH NOTE: this template sits at resume/applications/template.tex, beside
%%  template.cls, so the class path is just `template`. On copy into
%%  applications/{year}/{company}/{role}/, the orchestrator rewrites it to
%%  ../../../template (three levels up to applications/, where template.cls lives).
%% ============================================================================

\documentclass[
	11pt
]{../../../template}

\name{Ankur Desai}
\contact{(248) $\hspace{-0.1em}$ 657 $\hspace{-0.4em}$ -- $\hspace{-0.4em}$ 3805 \\ ardusa05@gmail.com \\ \href{https://ardusa.dev/}{www.ardusa.dev} \\ \href{https://linkedin.com/in/ardusa/}{linkedin.com/in/ardusa}}

\begin{document}

	% ---- EDUCATION (fixed; writer does not modify) --------------------------
	\begin{rSection}{E}{ducation}
		\begin{rSectionEntry}
			{Michigan State University}
			{Expected May 2028}
			{B.S. in Computer Science}
			{}
			{East Lansing, MI}

			\item \textbf{GPA:} 3.5 / 4.0
			\item \textbf{Coursework:} Data Structures \& Algorithms, OOP, Computer Architecture, Linear Algebra, Discrete Math
		\end{rSectionEntry}
	\end{rSection}

	% ---- EXPERIENCE (writer fills from context.md pool) ---------------------
	\begin{rSection}{E}{xperience}
		\begin{rSectionEntry}
			{Claude Builder Club @ MSU}
			{Oct 2025 -- Present}
			{Co-Founder \& Vice President}
			{}
			{East Lansing, MI}

			\item Built the club's mentorship lineage as an interactive 3D radial graph in react-three-fiber, hand-writing a proportional-arc layout that prevents edge crossings and a custom camera that autofits each family to the viewport.
			\item Co-founded and grew an MSU engineering community from zero to 150+ members, replacing a sprawl of spreadsheets, Slack threads, and manual invites with a single club-operations platform I designed and built.
			\item Tuned React Query cache tiers by data volatility across the dashboards, caching stable roles longer than fast-changing application queues to cut redundant refetches while keeping live data fresh.
			\item Automated member onboarding so an accepted applicant is provisioned into the right GitHub teams, Slack channels, Discord roles, and Google Drive access in a single flow, cutting onboarding from 15 minutes to seconds.
			\item Shipped the platform continuously from main on a GitHub Actions pipeline that applies database migrations and deploys edge functions on every push, taking the team from manual changes to reviewed continuous deploys.
		\end{rSectionEntry}

		\begin{rSectionEntry}
			{MindMosaic}
			{Aug 2025 -- May 2026}
			{Software Engineering Intern}
			{}
			{East Lansing, MI}

			\item Built and shipped an interactive knowledge-graph platform that keeps dense, interconnected content navigable as it grows, delivering a production MVP to real users on a polyglot Java Spring Boot, Neo4j, and PostgreSQL stack.
		\end{rSectionEntry}
	\end{rSection}

	% ---- PROJECTS (writer fills from context.md pool) -----------------------
	\begin{rSection}{P}{rojects}
		\begin{rSectionEntry}
			{Dadei}
			{\href{https://dadei.app}{dadei.app}}
			{Ambient voice assistant that turns overheard conversation into human-approved calendar, email, and task actions}
			{React, WebSocket, Web Audio}
			{}

			\item Built a React client that streams microphone audio up over a WebSocket with client-side Web Audio processing, receiving incremental live captions back in under 300ms as the backend transcribes each utterance.
			\item Routed all model-proposed side effects through one approval chokepoint with a cancelable countdown, giving every irreversible action one auditable seam between the LLM and the outside world.
		\end{rSectionEntry}

		\begin{rSectionEntry}
			{WizViz}
			{\href{https://devpost.com/software/wizviz}{devpost.com/software/wizviz}}
			{Webcam fighting game driven by real-time body-gesture recognition, winner of the Interactive Media track at SpartaHack X}
			{}
			{}

			\item Ran pose inference asynchronously off the render loop with a monotonic timestamp guard that drops out-of-order results, sustaining 60 FPS with sub-20ms input lag instead of blocking on per-frame detection.
			\item Classified five distinct combat gestures from body landmarks with a scale-invariant geometric model normalized by torso length, making detection work across players and camera distances without any labeled training data.
			\item Supported both local two-player, splitting players by landmark position across the frame with zero extra cameras, and single-player against a finite-state opponent that rests, defends, heals, or finishes by game state.
		\end{rSectionEntry}

		\begin{rSectionEntry}
			{fliks}
			{\href{https://github.com/fliks-gg}{github.com/fliks-gg}}
			{Multi-game platform that crowd-validates gameplay skill and mints shareable certification cards}
			{}
			{}

			\item Made ratings immutable and gated commenting behind a submitted rating, so the certification signal cannot be retroactively gamed and social participation feeds skill data instead of diluting it.
		\end{rSectionEntry}
	\end{rSection}

	% ---- TECHNICAL SKILLS (writer picks 3-5 buckets from context.md) --------
	\begin{rSection}{T}{echnical Skills}
		\begin{rSet}{Languages}{JavaScript/TypeScript, Python, SQL}
		\end{rSet}
		\begin{rSet}{Frameworks}{React, Spring Boot}
		\end{rSet}
		\begin{rSet}{Databases}{PostgreSQL, Neo4j}
		\end{rSet}
		\begin{rSet}{Tools}{Git, Figma, Linux}
		\end{rSet}
		\begin{rSet}{Cloud \& Infrastructure}{GitHub Actions}
		\end{rSet}
	\end{rSection}

\end{document}
